\documentclass[11pt]{amsart}
\usepackage[margin=1.1in]{geometry}
\usepackage[T1]{fontenc}
\usepackage{amsmath,amssymb,amsthm,mathtools,booktabs,hyperref}
\hypersetup{colorlinks=true,linkcolor=blue!50!black,citecolor=blue!50!black,urlcolor=blue!50!black}
\usepackage[dvipsnames]{xcolor}
\numberwithin{equation}{section}
\emergencystretch=4em
\newtheorem{theorem}{Theorem}[section]
\newtheorem{lemma}[theorem]{Lemma}
\newtheorem{proposition}[theorem]{Proposition}
\newtheorem{corollary}[theorem]{Corollary}
\theoremstyle{definition}
\newtheorem{definition}[theorem]{Definition}
\theoremstyle{remark}
\newtheorem{remark}[theorem]{Remark}
\newcommand{\li}{\operatorname{li}}
\newcommand{\Rp}{\operatorname{Re}}
\newcommand{\Ip}{\operatorname{Im}}
\newcommand{\eps}{\varepsilon}
\newcommand{\code}[1]{\texttt{#1}}
\DeclareMathOperator{\sgn}{sgn}

\title[Ramanujan's prime-counting inequality]{Ramanujan's prime-counting inequality holds for all $x \geq 38\,358\,837\,683$}
\author{Jaideep Sai Padhi}
\address{Purdue University, West Lafayette, IN, USA}
\email{jpadhi@purdue.edu}
\date{\today}
\subjclass[2020]{Primary 11N05; Secondary 11M26, 11Y35}
\keywords{Ramanujan's inequality, prime-counting function, zero-density estimates, zero-free regions, explicit formula, interval arithmetic}

\begin{document}
\begin{abstract}
Ramanujan observed that $\pi(x)^2 < \frac{e x}{\log x}\,\pi\!\left(\frac{x}{e}\right)$ for all sufficiently large $x$.
Dudek and Platt showed, assuming the Riemann hypothesis, that the last counterexample is $x = 38\,358\,837\,682$.
Unconditionally the inequality was known for $38\,358\,837\,683 \le x \le e^{103}$ (Johnston) and for $x \ge e^{3158.442}$ (Axler).
We close the gap: the inequality holds unconditionally for every $x \ge 38\,358\,837\,683$.
The proof uses no verification of the Riemann hypothesis beyond the height $3\cdot 10^{12}$ established by Platt and Trudgian.
It combines a smoothed explicit formula of B\"uthe's type with three new explicit inputs:
\begin{itemize}
\item a log-free zero-density estimate near the $1$-line, obtained by reworking Bellotti's argument with a sharp asymptotic evaluation of the Barban--Vehov mean square;
\item a height-dependent zero-free region, obtained by evaluating the inequality of Bellotti, Trudgian and Yang exactly;
\item a repulsion lemma showing that two hypothetical zeros close to the $1$-line cannot lie close to each other in height.
\end{itemize}
Every numerical step is certified in rigorous interval (ball) arithmetic.
Along the way we record several corrections to the explicit inputs we use, none of which affects the published main results.
\end{abstract}
\maketitle
\tableofcontents
